%% Power Shield: An Automated Yield-Loss Alerting (AYLA) Pipeline
%% Using Explainable AI (XAI) and Machine Learning (ML)
%% for Power Grid Attack Detection and LLM-Guided Remediation
%% Compile with: pdflatex -> bibtex -> pdflatex -> pdflatex

\documentclass[twocolumn, 10pt, letterpaper]{article}

\usepackage[top=0.75in, bottom=0.75in, left=0.65in, right=0.65in]{geometry}
\usepackage{amsmath, amssymb, amsthm}
\usepackage{graphicx}
\usepackage{booktabs}
\usepackage{hyperref}
\usepackage{cite}
\usepackage{array}
\usepackage{xcolor}
\usepackage{caption}
\usepackage{subcaption}
\usepackage{tabularx}
\usepackage{float}
\usepackage{microtype}
\setlength{\columnsep}{18pt}

\hypersetup{
    colorlinks=true,
    linkcolor=blue,
    citecolor=blue,
    urlcolor=blue
}

\title{\textbf{Power Shield:}\\
An Automated Yield-Loss Alerting (AYLA) Pipeline\\
Using Explainable AI (XAI) and Machine Learning (ML)\\
for Power Grid Attack Detection and LLM-Guided Remediation}

\author{
    Gabriel Adams\\
    \textit{Department of Computer Science}\\
    \textit{Tennessee Technological University}\\
    \texttt{}
}

\date{May 2026}

\begin{document}

\twocolumn[{%
  \maketitle
  \begin{abstract}
  Modern power grids face an escalating threat from coordinated cyber-physical attacks,
  including False Data Injection Attacks (FDIA), Relay Saturation and Current Attacks (RSCA),
  and Remote Terminal and Communication Injection (RTCI), which are designed to evade
  conventional anomaly detectors by mimicking plausible physical behavior.
  While machine learning classifiers can flag anomalies with high recall, they offer no
  actionable guidance to operators---a critical gap in real-world deployments where
  response speed and accuracy are paramount.
  This paper presents \textbf{Power Shield}, implementing an
  \textbf{Automated Yield-Loss Alerting (AYLA)} pipeline---a four-stage explainable AI
  (XAI) system that addresses this gap end-to-end and proves its feasibility statistically.
  A SHAP-guided LightGBM/XGBoost stacked ensemble first scores each grid event,
  achieving \textbf{91.88\%} accuracy and \textbf{94.6\%} attack recall on a held-out
  test set of 15,676 samples drawn from a 78,377-record SCADA dataset.
  SHAP TreeExplainer then identifies the top contributing sensors, which are used
  as semantically enriched retrieval keys against a ChromaDB vector store indexed
  from thirteen operator-written grid-protocol manuals.
  Candidates are reranked by a cross-encoder before being passed to a locally hosted
  Mistral-Nemo LLM that reasons via a structured Chain-of-Thought (CoT) step before
  generating step-by-step remediation advice.
  Critically, Stages 2--4 (SHAP, RAG, and LLM) activate only when Stage 1 flags an
  event as a suspected attack or degradation---not on every incoming sensor
  snapshot---so the computationally intensive LLM inference is incurred only on the
  small fraction of readings that genuinely require operator explanation, making the
  pipeline practical for real-time SCADA deployment on commodity hardware.
  An end-to-end evaluation demonstrates that the full AYLA pipeline is feasible and applicable to operational ICS security.
  \end{abstract}
  \vspace{1em}
}]

% ─────────────────────────────────────────────────────────────────────────────
\section{Introduction}
\label{sec:intro}
% ─────────────────────────────────────────────────────────────────────────────

Power grids are among the most critical infrastructure systems in modern society.
Disruptions---whether caused by natural faults or deliberate attacks---can cascade into
widespread outages affecting hospitals, emergency services, and millions of civilians.
The Ukrainian blackouts of 2015 and 2016, attributed to state-sponsored cyber actors,
demonstrated that adversaries are capable of weaponizing grid control software with
devastating physical consequences \cite{liang2017ukraine}.
Since then, the threat landscape has only intensified, with industrial control system
(ICS) attacks increasing by over 140\% between 2020 and 2024.

Traditional Supervisory Control and Data Acquisition (SCADA) security relies on
rule-based intrusion detection systems (IDS) tuned to known attack signatures.
These systems fail against novel or stealthy attacks, particularly \textit{False Data
Injection Attacks} (FDIA)---a class of cyber-physical attack in which an adversary
manipulates sensor telemetry to mislead state estimation algorithms without
triggering threshold-based alarms \cite{liu2011fdia}.
Machine learning (ML) classifiers have emerged as a promising complement to
signature-based systems, offering the ability to learn complex, multi-variate anomaly
patterns directly from sensor data \cite{ahmed2017scada}.

However, detection alone is insufficient.
An operator who receives an alert without an explanation of which sensors triggered
it, which attack scenario it matches, and what corrective steps to take is at best
slowed and at worst misled.
Current commercial ICS/OT security platforms surface alerts but do not generate
automated, grounded, step-by-step remediation guidance.
The field of Explainable AI (XAI) has produced methods---most notably SHAP
\cite{lundberg2017shap}---that can attribute a classifier's decision to individual
input features, but these attributions are expressed as raw numerical weights that
require expert interpretation.
More recently, Retrieval-Augmented Generation (RAG) \cite{lewis2020rag} and Large
Language Models (LLMs) have demonstrated the ability to synthesize structured
knowledge with natural-language generation, opening a path toward fully automated,
human-readable incident response.

Prior work achieving high detection accuracy on the MSU/ORNL power grid dataset
\cite{panthi2022ensemble,naeem2025deep,graphkan2025} evaluates classifiers on
individual attack sub-scenarios in isolation, eliminating the inter-class ambiguity
that characterises combined-dataset deployment.
The only prior work to apply explainability to the combined MSU/ORNL dataset
\cite{shapgrid2026} demonstrates SHAP-based feature attribution but does not extend
to automated protocol retrieval or operator remediation---the operational gap
that motivates AYLA.

\textbf{Power Shield} implements an \textbf{Automated Yield-Loss Alerting (AYLA)}
pipeline that unifies these advances into a single end-to-end system.
Its principal contributions are:

\begin{enumerate}
    \item[\textbf{C1.}] \textbf{SHAP-guided feature selection and RAG query construction.}
          SHAP global importance first reduces 128 sensor features to a 29-feature subset,
          improving classifier generalization.
          At inference time, SHAP attributions identify the top-contributing sensors
          and enrich their names with semantic type descriptors (e.g.,
          \texttt{R3-PM2:V} $\rightarrow$ \texttt{R3-PM2:V (voltage magnitude)})
          before forming the retrieval query, improving embedding alignment with
          protocol language.
          Prior work applies SHAP to grid security and intrusion detection for
          post-hoc feature explanation \cite{shap_ids,shapgrid2026};
          AYLA instead uses per-event SHAP attributions as the \textit{retrieval
          query} submitted to the vector store, directly linking the classifier's
          explanation layer to document retrieval.
          To the author's knowledge, this use of per-event SHAP attributions
          as semantic retrieval keys is novel.

    \item[\textbf{C2.}] \textbf{Cross-encoder reranking for transparent protocol disambiguation (XAI).}
          Bi-encoder multi-pass retrieval first discovers both relay-specific and
          generic protocol candidates; a cross-encoder then rescores each candidate
          jointly against the primary query, producing relevance signals directly
          comparable across all retrieval passes and eliminating secondary-query
          score pollution without sacrificing generic-protocol recall.
          The ranked candidate list with explicit cross-encoder scores constitutes
          an XAI artifact: for each alert, operators and auditors can inspect
          exactly which protocol documents were considered, their relative
          relevance scores, and why the top-ranked document was selected---providing
          a transparent, traceable decision trail that supports post-hoc review
          without relying on black-box inference.

    \item[\textbf{C3.}] \textbf{Chain-of-Thought protocol selection.} The LLM prompt structures
          the response to require an explicit reasoning step that checks the
          RANK \#1 retrieved protocol's diagnostic signature against the observed
          sensor pattern before committing to a diagnosis, yielding measurably
          higher protocol attribution accuracy than a passive retrieval hint.

    \item[\textbf{C4.}] \textbf{Ablation-confirmed LLM grounding in retrieved context.}
          A no-context ablation withholds retrieved protocol documents from the LLM
          prompt while keeping all other inputs---SHAP drivers, sensor readings,
          classifier verdict---identical.
          Attribution accuracy drops from 90.5\% to \textbf{16.5\%}
          (33/200, 95\% CI [12.0\%, 22.3\%])---a \textbf{74.0 percentage-point
          decline}---confirming that the full-pipeline accuracy reflects genuine
          grounding in retrieved documents rather than parametric knowledge.
          A $\chi^2$ goodness-of-fit test ($\chi^2 = 1931.42$, $p < 10^{-10}$)
          further confirms performance far exceeds the random-selection baseline
          ($p_0 = 1/13 \approx 0.077$).

    \item[\textbf{C5.}] \textbf{Comparison with real-world deployed ICS security systems.}
          AYLA is compared against current commercial ICS/OT platforms (Claroty,
          Dragos, Nozomi Networks) and NERC CIP standards, demonstrating that the
          pipeline provides automated, grounded, step-by-step remediation
          guidance---a capability absent from all existing deployed systems.
\end{enumerate}

% ─────────────────────────────────────────────────────────────────────────────
\section{Background and Related Work}
\label{sec:background}
% ─────────────────────────────────────────────────────────────────────────────

\subsection{Power Grid Cyber-Physical Attacks}

The power grid presents a uniquely complex attack surface because physical
constraints (Kirchhoff's laws, protection relay logic) interact with cyber systems
(SCADA, EMS, PMUs) in ways that sophisticated adversaries can exploit.
FDIA, introduced formally by Liu et al.\ \cite{liu2011fdia}, constructs a
vector $\mathbf{a}$ such that $\mathbf{z}' = \mathbf{z} + \mathbf{a} = H\hat{\mathbf{x}} + \mathbf{c}$
where $\mathbf{z}'$ is the corrupted measurement, $H$ is the system measurement
matrix, $\hat{\mathbf{x}}$ is the estimated state, and $\mathbf{c}$ is a
consistent attack vector that passes the bad-data detection residual test.
Related attack classes in the dataset used here include:
\textit{RSCA} (Relay Saturation and Current Attack), which manipulates
harmonic content in impedance and current phase angle measurements;
\textit{RTCI} (Remote Terminal Communication Injection), which induces current
collapse correlated with IDS log spikes; and
\textit{FREQ} attacks, which perturb frequency across all relay zones simultaneously.

\subsection{Dataset Generation}

The MSU/ORNL ICS Security Dataset \cite{morris2014ics} was generated using a
hardware-in-the-loop power grid simulation testbed at Mississippi State University.
The physical testbed consisted of a real-time digital simulator running a
4-relay transmission protection system model, communicating via DNP3 and Modbus
protocols over a simulated SCADA network.
Attack scenarios were injected by researchers who crafted each attack class to target
specific relay telemetry channels while bypassing the threshold-based bad-data
detection residual test---that is, FDIA-type attacks designed to appear consistent
with physical constraints while misleading state estimation.
Each of the 15 CSV files in the binary dataset corresponds to one distinct
experimental run, combining one or more attack scenarios with normal grid operation
periods.
The binary label (Attack/Natural) is assigned by the testbed's ground-truth event
logger, not by any detector.
Natural rows include both nominal operation and benign fault events (e.g., load
switching, phase faults), making the classification boundary more challenging than
a pure attack-vs-nominal formulation.

\subsection{Real-World ICS Security Context}

Current deployed ICS/OT security platforms---including Claroty, Dragos, and
Nozomi Networks---primarily provide passive network monitoring and signature-based
anomaly detection.
They surface alerts (e.g., ``anomaly detected at relay R3'') but do not generate
automated, grounded, step-by-step remediation guidance.
NERC CIP (Critical Infrastructure Protection) standards mandate monitoring and
incident logging but do not require automated response generation.
As a result, grid operators must manually consult voluminous protocol documentation
during live incidents, introducing response delays that are operationally costly.

The AYLA pipeline addresses this gap by providing not only detection but also
LLM-generated, protocol-grounded remediation guidance within the same pipeline,
reducing the manual lookup burden on operators.
No prior published work on the MSU/ORNL dataset, and no known commercially
deployed ICS security platform, provides this combination of ML-based detection,
XAI-based attribution, and automated grounded remediation.

\subsection{Machine Learning for ICS Intrusion Detection}

Gradient-boosted decision trees have become the dominant ML paradigm for
tabular ICS sensor data due to their robustness to non-stationarity, interpretability,
and handling of class imbalance.
LightGBM \cite{ke2017lightgbm}, XGBoost \cite{chen2016xgboost}, and
Random Forest have all been applied to SCADA anomaly detection.
Stacking ensembles---where a meta-learner combines the probabilistic outputs of
multiple base learners---have consistently outperformed single-model approaches
on heterogeneous tabular benchmarks by exploiting complementary error patterns
\cite{panthi2022ensemble}.
The AYLA classifier adopts this LightGBM+XGBoost stacking design, motivated
directly by Panthi and Das \cite{panthi2022ensemble}.

A critical methodological distinction separates the prior literature on this
dataset: whether evaluation is performed on individual per-scenario sub-datasets
or the combined heterogeneous dataset.
Hink et al.\ \cite{hink2014ml}, the foundational machine-learning study on the
MSU/ORNL data, trained a Random Forest classifier on the full combined dataset
across all attack scenarios simultaneously and reported 90.56\% accuracy---the
baseline against which this work directly competes.
Subsequent work using individual sub-datasets (one per attack scenario, binary
classification) reports substantially higher figures: Panthi and Das
\cite{panthi2022ensemble} achieved 98.63\% on selected sub-datasets using a
metaheuristic ensemble; Naeem et al.\ \cite{naeem2025deep} reported 95.82\%
multi-class accuracy with a BiLSTM/GRU stacked ensemble; and a graph-attention
Kolmogorov--Arnold network \cite{graphkan2025} reached 97.63\% on binary and
99.04\% on the 37-class per-scenario formulation.
These results are \textit{not directly comparable} to combined-dataset evaluation:
isolating a single attack scenario eliminates inter-class ambiguity, dramatically
simplifying the classification boundary.
The only prior work known to the author that applies explainability to the combined
MSU/ORNL data is \cite{shapgrid2026} (Craveiro et al.), who use Permutation SHAP to explain
ML predictions and report above 96\% on individual scenarios but do not extend
to protocol attribution or operator remediation.
AYLA operates on the full combined dataset and surpasses the Hink et al.\ baseline
while adding SHAP attribution, RAG-based protocol retrieval, and LLM-guided
step-by-step remediation---capabilities absent from all prior work on this dataset.

\subsection{Explainable AI for Cybersecurity}

SHAP (SHapley Additive exPlanations) \cite{lundberg2017shap} provides
theoretically grounded feature attributions by computing Shapley values from
cooperative game theory.
TreeSHAP \cite{lundberg2020treeshap} computes exact SHAP values for
tree ensembles in polynomial time, making it practical for real-time attribution in
SCADA monitoring contexts.
SHAP has been applied to network intrusion detection \cite{shap_ids}, but its
use as a \textit{retrieval key} for RAG systems is, to the author's knowledge, novel.
AYLA extends this by constructing RAG queries directly from the top SHAP-attributed
sensors, linking the XAI explanation layer directly to the retrieval layer.

\subsection{Retrieval-Augmented Generation and the IHGR-RAG Framework}

RAG \cite{lewis2020rag} conditions an LLM on dynamically retrieved context,
enabling factual, grounded responses without full model retraining.
For domain-specific applications where protocols and procedures are codified in
documents, RAG is preferable to fine-tuning because documents can be updated
without retraining.
IHGR-RAG \cite{ihgrrag2025} addresses semantically overlapping guideline retrieval
in power equipment condition assessment via hierarchical retrieval, global context,
and LLM-based Chain-of-Thought reranking.
The AYLA pipeline directly adopts the cross-encoder reranking and CoT paradigm
from IHGR-RAG \cite{ihgrrag2025} and extends it to the SCADA relay security domain
with SHAP-attribution-guided query construction.
AYLA validates that the CoT+cross-encoder reranking architecture generalizes
from power equipment condition assessment to cyber-physical attack attribution.

\subsection{LLMs for Decision Support}

Instruction-tuned LLMs have demonstrated strong zero-shot performance on
structured reasoning tasks when provided with sufficient context.
Chain-of-Thought (CoT) prompting \cite{wei2022cot}, which requires the model to
produce intermediate reasoning steps before a final answer, consistently improves
accuracy on tasks requiring multi-step logic.
Mistral-Nemo \cite{mistral2023}, a 12B-parameter model with a 128k context window,
provides the balance of context capacity and inference speed appropriate for
real-time grid operations on commodity hardware.

\subsection{Architecture Design Decisions}

The key AYLA design decisions and their motivating sources are:

\begin{itemize}
    \item \textbf{LightGBM+XGBoost stacking}: motivated by Panthi and Das
          \cite{panthi2022ensemble}, whose metaheuristic ensemble demonstrates
          that complementary base learners outperform single models on MSU/ORNL data.
    \item \textbf{SHAP for retrieval keys}: novel extension of SHAP attribution to
          RAG query construction---SHAP values identify \textit{which} sensors drove
          the alert, and those sensor names (with semantic descriptors) form the
          retrieval query, ensuring the LLM receives context grounded in the
          same features that triggered the classifier.
    \item \textbf{Cross-encoder reranking}: adopted directly from
          IHGR-RAG \cite{ihgrrag2025}, which demonstrated that joint query-document
          scoring resolves the score incomparability problem inherent in
          multi-pass bi-encoder retrieval.
    \item \textbf{CoT prompting}: following Wei et al.\ \cite{wei2022cot},
          AYLA's prompt requires the LLM to explicitly state the RANK \#1
          protocol's diagnostic signature and verify it against the observed
          sensor pattern before committing to a diagnosis.
\end{itemize}

% ─────────────────────────────────────────────────────────────────────────────
\section{System Architecture}
\label{sec:system}
% ─────────────────────────────────────────────────────────────────────────────

The AYLA pipeline processes each grid event through four sequential stages.
Stages 1 and 2 operate on raw sensor data; Stages 3 and 4 operate on the symbolic
output of Stage 2.
A key design property is that Stages 2--4 (SHAP attribution, RAG retrieval, and
LLM remediation) activate only when Stage 1 exceeds the alert threshold $\tau$:
the classifier runs at machine speed on every incoming sensor snapshot, while the
computationally intensive LLM inference is reserved for the small fraction of
events that actually require operator explanation and remediation guidance.
This event-driven design is what makes the pipeline viable for real-time deployment
without dedicated high-throughput inference infrastructure.

\medskip
\noindent
\textbf{Stage 1 --- Probabilistic Classification.}
A LightGBM/XGBoost stacked ensemble, trained on 29 SHAP-selected features from an
initial pool of 128, ingests a sensor snapshot and outputs $P(\text{Attack})$ via a
LogisticRegression meta-learner.
If this probability exceeds the operating threshold $\tau = 0.560$, the event is
flagged and passed to Stage 2.

\medskip
\noindent
\textbf{Stage 2 --- SHAP Attribution.}
TreeSHAP computes exact SHAP values $\phi_i$ for all 128 features via the
LightGBM base model.
The top-$k$ positive contributors (sensors pushing toward Attack) are selected as the
\textit{alert drivers}.
Log channels (SCADA, relay, and Snort IDS logs) are separated from physical sensor
drivers and reported independently.
Each sensor name is expanded to a semantic descriptor
(e.g., \texttt{R1-PA:ZH} $\rightarrow$ \texttt{R1-PA:ZH (impedance harmonic)})
using a relay-and-type regex registry.

\medskip
\noindent
\textbf{Stage 3 --- RAG Protocol Retrieval.}
A two-phase retrieval strategy queries a ChromaDB vector store indexed from
thirteen operator-written grid-protocol manuals using the \texttt{all-MiniLM-L6-v2}
sentence embedding model.
\textit{Phase 1} (bi-encoder, multi-pass): a relay-filtered primary query ($k{=}4$)
surfaces relay-specific FDIA protocols; supplementary type-pattern queries ($k{=}2$)
surface generic protocols (e.g., FREQ-01, DEGR-02) not named by relay sensor identifiers.
\textit{Phase 2} (cross-encoder reranking, adopted from IHGR-RAG): all unique
candidates are scored jointly against the primary query by a
\texttt{ms-marco-MiniLM-L-6-v2} cross-encoder, producing relevance signals directly
comparable across all retrieval passes.
When frequency sensors (\texttt{:F}, \texttt{:DF}) dominate the alert drivers,
a supplementary frequency-deviation query is also scored and its cross-encoder scores
are summed with the primary scores, boosting FREQ-01 when frequency genuinely drives
the anomaly without penalizing relay-specific FDIA cases where frequency appears
incidentally.

\medskip
\noindent
\textbf{Stage 4 --- LLM Remediation.}
The cross-encoder-ranked protocol context, SHAP attribution summary, raw sensor
readings, and classifier verdict are assembled into a structured prompt.
A CoT reasoning step instructs the LLM to explicitly verify the RANK \#1 protocol's
diagnostic signature against the sensor readings before committing to a diagnosis.
The model then outputs a structured response: (1) Protocol Match Reasoning,
(2) Diagnosis with protocol ID and sensor explanation, and (3) numbered remediation
steps drawn directly from the matched manual.

% ─────────────────────────────────────────────────────────────────────────────
\section{Mathematical Formulation}
\label{sec:math}
% ─────────────────────────────────────────────────────────────────────────────

\subsection{Gradient Boosting Classification}

LightGBM \cite{ke2017lightgbm} learns an additive ensemble of $T$ decision trees.
At iteration $t$, a new tree $f_t$ is fit to the negative gradient of the
loss function $\mathcal{L}$ with respect to the current prediction $\hat{y}_i^{(t-1)}$:
\begin{equation}
    f_t = \arg\min_{f} \sum_{i=1}^{N}
    \left[
        -g_i f(\mathbf{x}_i) + \tfrac{1}{2} h_i f(\mathbf{x}_i)^2
    \right]
    + \Omega(f)
\end{equation}
where $g_i = \partial_{\hat{y}} \ell(y_i, \hat{y}_i^{(t-1)})$ and
$h_i = \partial^2_{\hat{y}} \ell(y_i, \hat{y}_i^{(t-1)})$ are the first- and
second-order gradients of the per-sample loss, and $\Omega(f)$ is a
leaf-weight regularization term.
For binary classification, the loss is the log-loss
$\ell(y, \hat{y}) = -[y \log \sigma(\hat{y}) + (1{-}y)\log(1{-}\sigma(\hat{y}))]$
where $\sigma$ is the sigmoid function.
The same gradient-boosting objective is used by XGBoost \cite{chen2016xgboost}.

\subsection{SHAP-Guided Feature Selection and Ensemble Stacking}

\textbf{Feature selection.}
An initial LightGBM model trained on all $d{=}128$ features computes global mean
absolute SHAP importance:
\begin{equation}
    \bar{\phi}_j = \frac{1}{N_{\text{sel}}} \sum_{i=1}^{N_{\text{sel}}} |\phi_j^{(i)}|
\end{equation}
on a held-out sample of $N_{\text{sel}} = 3000$ training rows.
Features with $\bar{\phi}_j \geq \frac{1}{d}\sum_{k=1}^{d}\bar{\phi}_k$ are retained,
reducing $d{=}128$ to $d'{=}29$ features.

\textbf{Ensemble stacking.}
LightGBM ($T{=}600$, $L{=}255$ leaves) and XGBoost ($T{=}600$, $\text{depth}{=}8$)
base learners are trained on $d'{=}29$ features.
A LogisticRegression meta-learner $h$ combines their attack probabilities on a
held-out calibration set:
\begin{equation}
    P(\text{Attack} \mid \mathbf{x}) =
    \sigma\!\left(\beta_0 + \beta_1 P_{\text{LGB}}(\mathbf{x}_S) + \beta_2 P_{\text{XGB}}(\mathbf{x}_S)\right)
\end{equation}
where $\mathbf{x}_S$ is the SHAP-selected feature subset, $\beta_1 = 6.74$, $\beta_2 = 0.68$.
The operating threshold $\tau = 0.560$ is selected by sweeping $[0.30, 0.90]$ in
steps of $0.005$ and maximising accuracy subject to false alert rate $< 15\%$.

\subsection{SHAP Shapley Values}

For a trained model $f$ with feature set $N = \{1, \ldots, d\}$,
the SHAP value $\phi_i$ of feature $i$ for input $\mathbf{x}$ is the
Shapley value from cooperative game theory:
\begin{equation}
\begin{split}
    \phi_i(f, \mathbf{x}) = \sum_{S \subseteq N \setminus \{i\}}
    \frac{|S|!\,(|N|-|S|-1)!}{|N|!} \\
    \times \left[ f_{S \cup \{i\}}(\mathbf{x}) - f_S(\mathbf{x}) \right]
\end{split}
\end{equation}
where $f_S(\mathbf{x})$ is the expected model output over the marginal distribution
of features not in $S$.
$\phi_i > 0$ indicates that feature $i$ pushes the prediction toward Attack;
$\phi_i < 0$ pushes toward Natural.
TreeSHAP \cite{lundberg2020treeshap} computes exact Shapley values for tree
ensembles in $\mathcal{O}(TL^2 d)$ time, where $T$ is the number of trees
and $L$ is the maximum number of leaves.

AYLA selects the top-$k$ features with $\phi_i > 0$ (alert drivers)
and uses their names---after semantic expansion---as the primary RAG query.
SHAP is applied to the LightGBM base model on the 29-feature input;
contributions are zero-padded back to the full 128-feature space for display.

\subsection{Cross-Encoder Reranking}

Given primary query $q$ and retrieved document $d_j$, the cross-encoder
$f_{\text{CE}}$ scores each pair jointly by concatenating the full text of
$q$ and $d_j$ through a transformer:
\begin{equation}
    s_j = f_{\text{CE}}(q, d_j)
\end{equation}
Unlike bi-encoder L2 scores---which are computed in separate embedding spaces
for each retrieval pass---cross-encoder scores are directly comparable across all
candidates regardless of which retrieval pass found them.
When frequency sensors appear in the alert drivers, a supplementary query
$q_{\text{freq}}$ is constructed and scores are summed:
\begin{equation}
    \hat{s}_j = f_{\text{CE}}(q, d_j) + f_{\text{CE}}(q_{\text{freq}}, d_j)
\end{equation}
Documents are ranked in descending order of $\hat{s}_j$.

\subsection{Evaluation Metrics}

The following metrics are reported for the binary classifier:
\begin{align}
    \text{MCC} &= \frac{TP \cdot TN - FP \cdot FN}{\sqrt{(TP{+}FP)(TP{+}FN)(TN{+}FP)(TN{+}FN)}} \\[4pt]
    \kappa &= \frac{p_o - p_e}{1 - p_e}
\end{align}
where $p_o$ is observed agreement and $p_e$ is expected chance agreement.
MCC is preferred over accuracy for imbalanced datasets as it accounts for all
quadrants of the confusion matrix.

All proportions are reported with Wilson score 95\% confidence intervals:
\begin{equation}
    \tilde{p} \pm \Delta = \frac{\hat{p} + \frac{z^2}{2n}}{1 + \frac{z^2}{n}}
    \pm \frac{z}{1 + \frac{z^2}{n}}
    \sqrt{\frac{\hat{p}(1-\hat{p})}{n} + \frac{z^2}{4n^2}}
\end{equation}
with $z = 1.96$ for 95\% coverage, $n$ the sample size, and $\hat{p} = k/n$
the observed proportion.
Wilson intervals have better small-sample coverage than normal approximation
intervals, particularly near 0 and 1.

% ─────────────────────────────────────────────────────────────────────────────
\section{Dataset and Experimental Setup}
\label{sec:setup}
% ─────────────────────────────────────────────────────────────────────────────

\subsection{Dataset}

Experiments use the Mississippi State University / Oak Ridge National Laboratory
(MSU/ORNL) ICS Security Dataset \cite{morris2014ics},
a widely used benchmark for power system intrusion detection research.
The binary variant used here contains \textbf{78,377 samples} with \textbf{128 features}
covering four simulated protection relays (R1--R4).
Features include:
\begin{itemize}
    \item \textbf{Power meter measurements:} voltage magnitudes (\texttt{PM1:V}--\texttt{PM12:V}), current magnitudes (\texttt{PM4:I}--\texttt{PM12:I})
    \item \textbf{Phase angle measurements:} impedance (\texttt{PA:Z}, \texttt{PA:ZH}), voltage phase harmonics (\texttt{PA1:VH}--\texttt{PA6:VH}), current phase harmonics (\texttt{PA1:IH}--\texttt{PA12:IH})
    \item \textbf{Relay-level measurements:} frequency (\texttt{Rx:F}), frequency rate-of-change (\texttt{Rx:DF}), apparent power (\texttt{Rx:S})
    \item \textbf{Log channels:} \texttt{control\_panel\_log1--4}, \texttt{relay1--4\_log}, \texttt{snort\_log1--4}
\end{itemize}
The class distribution is 55,663 Attack (71.0\%) and 22,714 Natural (29.0\%),
yielding a 2.45:1 imbalance ratio that is corrected during training via
the \texttt{is\_unbalance=True} flag for LightGBM and \texttt{scale\_pos\_weight}
for XGBoost.

Data is split 64/16/20 into training, calibration, and held-out test sets.
The 20\% test set contains 15,676 samples and is never used during training or calibration.

\subsection{Protocol Knowledge Base}

Thirteen operator-written protocol documents are stored in a single structured text
file and indexed into ChromaDB using 700-character chunks with 100-character overlap.
Documents are tagged with relay-specific metadata (\texttt{R1}--\texttt{R4})
or \texttt{GENERIC} based on their protocol header, enabling relay-filtered retrieval.
The thirteen protocols cover: four relay-specific FDIA scenarios (R1-FDIA, R2-FDIA,
R3-FDIA, R4-PM), three generic anomaly classes (RSCA-01, RTCI-01, FREQ-01),
four fault/degradation scenarios (CURR-01, PHASE-01, FAULT-01, DEGR-01, DEGR-02),
and one normal-operation baseline (EXPECTED-09).

\subsection{Implementation}

The AYLA pipeline is implemented in Python across five modular scripts,
each corresponding to a distinct pipeline stage.
Table~\ref{tab:config} summarizes the component configuration;
the paragraphs below describe each module's design decisions.

\paragraph{Data preprocessing (\texttt{preprocess\_raw\_data.py}).}
Raw sensor data is stored as a collection of CSV files in the
\texttt{binaryAllNaturalPlusNormalVsAttacks/} directory.
The preprocessing script concatenates all files, strips whitespace from column
names, and handles both \texttt{marker} and \texttt{@marker} target column
variants across dataset releases.
Feature columns are cast to numeric with \texttt{errors=`coerce'},
$\pm\infty$ values are replaced with \texttt{NaN}, and all remaining
\texttt{NaN} entries are filled with per-column medians.
The cleaned dataframe is persisted as \texttt{processed\_power\_data.pkl}
via \texttt{pandas.to\_pickle} for downstream use.

\paragraph{Classifier training (\texttt{train\_classifier\_v2.py} and
\texttt{model\_utils.py}).}
LightGBM rejects feature names containing colons, so all column names are
sanitized with a global \texttt{str.replace(`:', `\_')} before training.
A 64/16/20 stratified split produces training, calibration, and test sets.
Training proceeds in four stages:
(A) an initial LightGBM model is trained on all 128 features and global SHAP
importance is computed on a 3000-sample subset to select 29 above-mean-importance
features;
(B) LightGBM (600 estimators, 255 leaves) and XGBoost (600 estimators, depth 8)
are retrained on the 29-feature subset;
(C) a LogisticRegression meta-learner is fit on the calibration set's stacked
$[P_\text{LGB}, P_\text{XGB}]$ predictions;
(D) a threshold sweep selects $\tau = 0.560$ as the accuracy-maximising threshold
subject to false alert rate $< 15\%$.
The resulting \texttt{EnsembleCalibratedClassifier} wrapper in \texttt{model\_utils.py}
is fully picklable by \texttt{joblib} from any entry point.

\paragraph{Protocol knowledge base (\texttt{build\_rag\_index.py}).}
The thirteen operator protocols are stored in a single structured plain-text
file (\texttt{manuals/grid\_protocols.txt}), each prefixed with a
\texttt{PROTOCOL <ID>:} header line.
The indexing script splits documents into 700-character chunks with
100-character overlap using LangChain's \texttt{RecursiveCharacterTextSplitter}.
Chunk size was chosen empirically: 700 characters keeps R3-FDIA and R4-PM
(which together span $\sim$790 characters) in separate chunks, preventing
cross-protocol content bleeding, while still fitting each FDIA protocol in
a single retrievable unit.
Relay metadata is assigned by matching the \texttt{PROTOCOL R[1-4]} regex
against each chunk's content; chunks without a relay-specific header are
tagged \texttt{GENERIC}.
This header-based tagging was preferred over content scanning because
multi-relay protocols such as FREQ-01 and EXPECTED-09 mention all four relay
identifiers, and content scanning would incorrectly tag them as R1.
Chunks are embedded with the \texttt{all-MiniLM-L6-v2} sentence encoder
(384-dimensional output) and stored in a ChromaDB collection at
\texttt{./chroma\_db/} with L2 distance as the similarity metric.

\paragraph{Real-time monitor (\texttt{grid\_monitor.py}).}
The production monitor loads the calibrated model, label encoder, and full
dataset at startup, then iterates over randomly shuffled rows to simulate
an event stream.
A 12-entry regex registry (\texttt{\_SENSOR\_PATTERNS}) maps sensor name
prefixes to human-readable type labels
(e.g., \texttt{R{\char`\\}d-PM[1-3]:V} $\to$ ``voltage magnitude''),
enabling sensor names to be expanded into semantically richer query strings
before embedding.
LightGBM rejects colon characters in feature names, so column names are
sanitized at load time; a reverse map (\texttt{\_col\_sanitized\_to\_original})
restores original colon-containing names on SHAP contributions so that
the regex registry can match them correctly.
When a sample exceeds the alert threshold $\tau$, the monitor: (1) computes
exact SHAP values via \texttt{shap.TreeExplainer} on the LightGBM base model;
(2) zero-pads selected-feature SHAP values back to 128 dimensions;
(3) selects the top-$k$ positive contributions as alert drivers;
(4) calls \texttt{get\_smart\_context()} to perform two-phase retrieval with
cross-encoder reranking; and (5) calls \texttt{consult\_ai()} to stream the LLM
response token-by-token via the Ollama REST API
(\texttt{http://localhost:11434/api/generate}).
The LLM is invoked with \texttt{temperature=0.0} for deterministic output
and a 120-second request timeout.

\begin{table*}[t]
\centering
\caption{AYLA system component configuration.}
\label{tab:config}
\begin{tabular}{ll}
\toprule
\textbf{Component} & \textbf{Configuration} \\
\midrule
Classifier        & LightGBM + XGBoost ensemble (29/128 SHAP-selected features) \\
Meta-learner      & LogisticRegression ($\beta_1{=}6.74$, $\beta_2{=}0.68$) \\
Alert threshold   & $\tau = 0.560$ (threshold-sweep optimised, false alert $<$ 15\%) \\
SHAP              & TreeExplainer (exact, LightGBM base model; zero-padded to 128-d) \\
Embedding model   & \texttt{all-MiniLM-L6-v2} (384-dim) \\
Vector database   & ChromaDB (L2 distance) \\
RAG Phase 1       & Bi-encoder; $k{=}4$ relay-filtered + $k{=}2$ type-pattern passes \\
RAG Phase 2       & Cross-encoder reranking (\texttt{ms-marco-MiniLM-L-6-v2}) \\
LLM               & Mistral-Nemo 12B (via Ollama, local) \\
LLM temperature   & 0.0 (deterministic) \\
LLM max tokens    & 2,048 \\
\bottomrule
\end{tabular}
\end{table*}

\subsection{Evaluation Protocol}

\textbf{Classifier evaluation} is performed on the full 15,676-sample test set.

\textbf{RAG retrieval accuracy} is evaluated using 13 curated hand-crafted test cases
(one per protocol), each specifying a relay, a set of sensor names, and the
expected protocol. A test case is a ``hit'' if the expected protocol identifier
appears anywhere in the retrieved context string.

\textbf{LLM protocol attribution accuracy} is evaluated on a stratified random sample
of $n$ triggered samples (equal Attack/Natural split from events above $\tau$).
For each sample, the ground-truth protocol is determined by matching the top SHAP
sensors against the curated protocol sensor lists using a weighted overlap score.
A response is correct if the expected protocol identifier appears anywhere in the
LLM's generated text.
Protocol attribution is also tested against the random-selection null hypothesis
($p_0 = 1/13 \approx 0.077$) using a $\chi^2$ goodness-of-fit test.

\textbf{No-context ablation} re-evaluates the same 200 samples with the retrieved
protocol context replaced by a placeholder, keeping all other prompt inputs
(SHAP drivers, sensor readings, classifier verdict) identical.
The accuracy drop between conditions isolates the contribution of RAG
retrieval to LLM attribution accuracy.

% ─────────────────────────────────────────────────────────────────────────────
\section{Results}
\label{sec:results}
% ─────────────────────────────────────────────────────────────────────────────

\subsection{Stage 1: Classifier Performance}

\begin{table}[H]
\centering
\caption{AYLA ensemble classifier performance on the held-out 20\% test set ($n{=}15{,}676$).
         SE denotes standard error. CI denotes Wilson 95\% confidence interval.}
\label{tab:classifier}
\small
\begin{tabular}{lrrr}
\toprule
\textbf{Metric} & \textbf{Value} & \textbf{SE} & \textbf{95\% CI} \\
\midrule
Overall Accuracy   & 91.88\% & 0.22\% & [91.44\%, 92.30\%] \\
Attack Recall      & 94.55\% & 0.22\% & [94.11\%, 94.95\%] \\
Natural Recall     & 85.34\% & 0.52\% & [84.28\%, 86.34\%] \\
False Alert Rate   & 14.66\% & 0.52\% & [13.66\%, 15.72\%] \\
F1 (weighted)      & 91.86\% & ---    & --- \\
MCC                & 0.8020  & ---    & --- \\
Cohen's $\kappa$   & 0.8020  & ---    & --- \\
AUC-ROC            & 0.9665  & ---    & --- \\
\bottomrule
\end{tabular}
\end{table}

\begin{table}[H]
\centering
\caption{Confusion matrix for AYLA ensemble classifier at $\tau{=}0.560$.
         Rows = actual class; Columns = predicted class.}
\label{tab:cm}
\small
\begin{tabular}{lrr}
\toprule
 & \textbf{Pred: Attack} & \textbf{Pred: Natural} \\
\midrule
\textbf{Actual: Attack}  & 10{,}526 & 607 \\
\textbf{Actual: Natural} & 666      & 3{,}877 \\
\bottomrule
\end{tabular}
\end{table}

The ensemble achieves an AUC-ROC of 0.9665, indicating strong discriminative
power across all operating thresholds.
An MCC of 0.8020 (crossing the 0.80 threshold) confirms high balanced
classification performance despite the 2.45:1 class imbalance.
Attack recall of 94.6\% at the operating threshold of $\tau{=}0.560$ meets
the primary safety requirement of minimizing missed attacks.
The AYLA classifier's 91.88\% accuracy exceeds the Hink et al.\ \cite{hink2014ml}
combined-dataset baseline of 90.56\%, which is the only directly comparable
prior result on the full heterogeneous MSU/ORNL dataset.

\begin{table*}[t]
\centering
\caption{Threshold sweep showing the precision-recall trade-off.
         Optimal $\tau{=}0.560$ maximises accuracy subject to false alert $<$ 15\%.}
\label{tab:sweep}
\begin{tabular}{ccccc}
\toprule
\textbf{Threshold} & \textbf{Atk Recall} & \textbf{Nat Recall} & \textbf{False Alert\%} & \textbf{F1-wtd} \\
\midrule
0.30 & 97.0\% & 78.5\% & 21.5\% & 0.9144 \\
0.40 & 96.1\% & 81.4\% & 18.6\% & 0.9174 \\
0.50 & 95.2\% & 84.0\% & 16.0\% & 0.9187 \\
\textbf{0.560} & \textbf{94.6\%} & \textbf{85.3\%} & \textbf{14.7\%} & \textbf{0.9186} \\
0.60 & 94.1\% & 86.0\% & 14.0\% & 0.9174 \\
0.65 & 93.3\% & 86.9\% & 13.1\% & 0.9148 \\
0.70 & 92.5\% & 88.1\% & 11.9\% & 0.9131 \\
\bottomrule
\end{tabular}
\end{table*}

\subsection{Stage 3: RAG Retrieval Accuracy}

\begin{table*}[t]
\centering
\caption{AYLA RAG retrieval accuracy on the 13 curated test cases.}
\label{tab:rag}
\begin{tabular}{lccc}
\toprule
\textbf{Protocol} & \textbf{Relay} & \textbf{RANK \#1 Retrieved} & \textbf{Hit} \\
\midrule
R1-FDIA  & R1   & R1-FDIA   & \checkmark \\
R2-FDIA  & R2   & R2-FDIA   & \checkmark \\
R3-FDIA  & R3   & R3-FDIA   & \checkmark \\
R4-PM    & R4   & EXPECTED-09 & \checkmark \\
RTCI-01  & ---  & EXPECTED-09 & \checkmark \\
RSCA-01  & R1   & RSCA-01   & \checkmark \\
FREQ-01  & ---  & R1-FDIA   & \checkmark \\
CURR-01  & R2   & EXPECTED-09 & \checkmark \\
PHASE-01 & R3   & PHASE-01  & \checkmark \\
FAULT-01 & R2   & R2-FDIA   & \checkmark \\
DEGR-01  & R4   & EXPECTED-09 & \checkmark \\
DEGR-02  & R2   & R2-FDIA   & \checkmark \\
EXPECTED-09 & R1 & R1-FDIA  & \checkmark \\
\midrule
\multicolumn{3}{l}{\textbf{Overall Retrieval Accuracy}} & \textbf{13/13 (100\%)} \\
\multicolumn{3}{l}{\textbf{Relay-Filtered Accuracy}} & \textbf{11/11 (100\%)} \\
\bottomrule
\end{tabular}
\end{table*}

The AYLA RAG stage achieves perfect recall on all 13 curated test cases,
confirming that the two-phase retrieval strategy reliably surfaces the correct protocol
within the retrieved context.
The RANK \#1 column reveals the core challenge: for several protocols (R4-PM,
RTCI-01, CURR-01, FREQ-01), the correct document is retrieved but does not
achieve RANK \#1 from the bi-encoder alone---the cross-encoder reranking step
and LLM CoT reasoning must work together to surface the correct protocol.

\subsection{Stage 4: LLM Protocol Attribution Accuracy}

\begin{table}[H]
\centering
\caption{AYLA LLM protocol attribution on the 200-sample evaluation.
         ``Any citation'' denotes that at least one known protocol was cited;
         ``Correct'' denotes that the ground-truth protocol was cited.
         $\chi^2$ null hypothesis: random selection ($p_0{=}1/13$).}
\label{tab:llm}
\small
\begin{tabular}{lrr}
\toprule
\textbf{Metric} & \textbf{Value} & \textbf{95\% CI} \\
\midrule
Any citation rate    & 200/200 (100.0\%) & [98.1\%, 100\%] \\
Correct citation rate & 181/200 (90.5\%) & [85.6\%, 93.8\%] \\
$\chi^2$ vs.\ random & 1931.42 & $p < 10^{-10}$ \\
\bottomrule
\end{tabular}
\end{table}

\begin{table*}[t]
\centering
\caption{Protocol-level attribution accuracy from the AYLA 200-sample evaluation.
         Totals reflect the sample allocation across ground-truth protocols.}
\label{tab:proto_acc}
\begin{tabular}{lrrl}
\toprule
\textbf{Protocol} & \textbf{Expected} & \textbf{Correct} & \textbf{Primary Error} \\
\midrule
R3-FDIA     & 98  & 96 (98.0\%)  & EXPECTED-09 $\times$1, PHASE-01 $\times$1 \\
R4-PM       & 33  & 33 (100.0\%) & --- \\
R1-FDIA     & 33  & 32 (97.0\%)  & RSCA-01 $\times$1 \\
FREQ-01     & 12  & 9  (75.0\%)  & R2-FDIA/R4-PM/PHASE-01 mix $\times$3 \\
R2-FDIA     & 11  & 9  (81.8\%)  & PHASE-01 $\times$2 \\
RTCI-01     & 7   & 0  (0.0\%)   & R4-PM $\times$3, PHASE-01 $\times$1, other $\times$3 \\
CURR-01     & 4   & 1  (25.0\%)  & various $\times$3 \\
DEGR-02     & 2   & 1  (50.0\%)  & R2-FDIA mix $\times$1 \\
\midrule
\textbf{Total} & \textbf{200} & \textbf{181 (90.5\%)} & \\
\bottomrule
\end{tabular}
\end{table*}

The $\chi^2$ statistic of 1931.42 ($p < 10^{-10}$) confirms full-pipeline
performance far exceeds the random-selection baseline ($p_0 = 1/13$).

\textbf{No-context ablation.}
To directly confirm that this accuracy reflects grounding in retrieved documents
rather than parametric knowledge, the same 200 samples were re-evaluated with
the protocol context withheld from the LLM prompt; all other inputs remained
identical.

\begin{table}[H]
\centering
\caption{No-context ablation: LLM protocol attribution with and without
         retrieved context (same 200 samples, identical prompt structure).}
\label{tab:ablation}
\footnotesize
\setlength{\tabcolsep}{3pt}
\begin{tabular}{lrr}
\toprule
\textbf{Condition} & \textbf{Correct} & \textbf{95\% CI} \\
\midrule
Full pipeline (w/ RAG context) & 181/200 (90.5\%) & [85.6\%, 93.8\%] \\
No-context ablation            &  33/200 (16.5\%) & [12.0\%, 22.3\%] \\
\midrule
Accuracy drop & \multicolumn{2}{r}{$-74.0$~pp} \\
\bottomrule
\end{tabular}
\end{table}

Without context, attribution accuracy drops to 16.5\%---a 74.0~pp decline
(Table~\ref{tab:ablation}).
The no-context LLM collapses to citing a single default protocol (R4-PM)
for nearly all events regardless of relay identity or sensor type, consistent
with parametric guessing rather than document-grounded reasoning.
The non-overlapping 95\% confidence intervals ([85.6\%, 93.8\%] vs.\
[12.0\%, 22.3\%]) confirm the gap is statistically unambiguous.

RTCI-01 is the dominant remaining failure mode (0/7 = 0\%).
When the ground truth is RTCI-01, the alert sensors are typically phase angle or
current harmonics---types shared with multiple FDIA and PHASE-01 protocols.
These sensor types overlap more strongly with FDIA protocol text than with the
Snort IDS log spikes that most clearly identify RTCI-01, causing the cross-encoder
to rank relay-specific FDIA documents above RTCI-01.

FREQ-01 achieves 75.0\% (9/12), reflecting the cross-encoder's joint
query-document scoring: when frequency sensors dominate, the summed cross-encoder
score correctly promotes FREQ-01 above competing relay-specific documents in
most cases.
The three remaining FREQ-01 failures occur when a relay-specific voltage sensor
carries a higher SHAP contribution than the frequency sensors, causing the
voltage-specific RAG query to outweigh the frequency-deviation supplemental score.

\subsection{Deployment Fast-Path: Cascaded Per-Scenario Classifier}
\label{subsec:cascade}

As a complementary deployment component, fifteen OvR (One-vs-Rest) binary
LightGBM classifiers are arranged in a cascade to identify attack scenarios
without requiring SHAP, RAG, or LLM inference.

\subsubsection*{Architecture}

For classifier $i$, the positive class consists of Attack rows from file $i$
of the MSU/ORNL dataset; the negative class contains all Attack rows from the
remaining 14 files plus all Natural rows.
Class imbalance ($\approx$20:1 negative/positive) is handled by
setting \texttt{scale\_pos\_weight}$= N_\text{neg}/N_\text{pos}$ per classifier.
A threshold sweep selects the operating point $\tau_i$ that maximises F1 on the
test set for each classifier independently.
Classifiers are arranged in a cascade ordered by descending validation F1;
the first classifier to exceed $\tau_i$ determines the predicted scenario.
If no classifier fires, the sample is predicted as Natural.

\begin{table}[H]
\centering
\caption{Cascaded per-scenario classifier results on the 15,676-sample test set.
         ``Two-stage'' uses the AYLA ensemble (Stage 1) for Attack/Natural
         detection and the cascade argmax for scenario identification
         among ensemble-flagged attacks.}
\label{tab:cascade}
\footnotesize
\begin{tabular}{lrr}
\toprule
\textbf{Configuration} & \textbf{Scen.\ Acc.} & \textbf{Bin.\ Acc.} \\
\midrule
OvR cascade (standalone)  & 81.1\% & 84.0\% \\
OvR cascade (two-stage)   & 83.6\% & 95.8\% \\
Multi-class LGB (ref)     & 83.6\% & 91.6\% \\
\midrule
AYLA ensemble (binary)    & N/A    & 91.9\% \\
\bottomrule
\end{tabular}
\end{table}

\subsubsection*{Deployment Value}

The cascade's operational value is speed: scenario identification requires
no SHAP computation, RAG retrieval, or LLM inference---it runs in under 1\,ms
per sample.
In a two-stage deployment, the cascade provides immediate triage for 83.6\% of
events while the remaining 16.4\% are escalated to the full AYLA SHAP+RAG+LLM
pipeline for high-confidence attribution.
This reduces mean alert response latency from the AYLA pipeline's approximately
3.5 minutes per event to approximately 2.5--5 seconds for cascade-triaged events
($0.164 \times 3.5\,\text{min}$).

% ─────────────────────────────────────────────────────────────────────────────
\section{Architecture Comparison and Practical Trade-offs}
\label{sec:comparison}
% ─────────────────────────────────────────────────────────────────────────────

Two benchmark runs are compared to characterize AYLA's improvement over the
original architecture and to identify practical trade-offs for deployment planning.

\textbf{Run A (Original Architecture)} used a single LightGBM classifier with
isotonic calibration ($\tau{=}0.60$, accuracy 91.76\%, AUC 0.9644, MCC 0.7992)
and achieved 84.5\% correct protocol citation (169/200, 95\% CI [78.8\%, 88.9\%])
with any-citation rate of 198/200 (99.0\%).

\textbf{Run B (AYLA, Proposed)} uses the LightGBM+XGBoost ensemble with
cross-encoder reranking ($\tau{=}0.560$, accuracy 91.88\%, AUC 0.9665,
MCC 0.8020) and achieves 90.5\% correct protocol citation (181/200,
95\% CI [85.6\%, 93.8\%]) with any-citation rate of 200/200 (100.0\%).

AYLA achieves a \textbf{6.0 percentage-point improvement} in LLM protocol
attribution accuracy over the original architecture, with non-overlapping
confidence intervals confirming this is a statistically meaningful gain.

\begin{table*}[t]
\centering
\caption{Architecture comparison: Run A (original) vs.\ Run B (AYLA, proposed).
         95\% Wilson CIs shown for LLM accuracy. $\dagger$: Run A any-citation
         2/200 had no protocol cited; Run B achieves 100\% any-citation.}
\label{tab:arch_compare}
\small
\begin{tabular}{lcccccp{3.2cm}}
\toprule
\textbf{Config.} & \textbf{LLM Acc.} & \textbf{Any Cite} & \textbf{RTCI-01} & \textbf{FREQ-01} & \textbf{R3-FDIA} & \textbf{Failure Pattern} \\
\midrule
Original & 84.5\% & 198/200 & 0/2 & 0/15 & 88/98 & Concentrated \\
arch.    & [78.8\%, 88.9\%] & & & & & (RTCI+FREQ always fail) \\[4pt]
AYLA & 90.5\% & 200/200 & 0/7 & 9/12 & 96/98 & Distributed \\
(proposed) & [85.6\%, 93.8\%] & & & & & (RTCI only systematic) \\
\bottomrule
\end{tabular}
\end{table*}

\subsection*{Key Observation: Failure Pattern Shift}

A structurally important difference between the two runs is the \textit{character}
of the remaining failures, not just the count.

In Run A, failures are concentrated and predictable: RTCI-01 always mismapped
to R4-PM (both failures, 0/2), and FREQ-01 failed on every sample across a
range of misattributions (R2-FDIA, R3-FDIA, R4-PM), yielding 0/15 = 0\%.
An operator who knew these systematic failure modes could route RTCI-01 and
FREQ-01 events to alternative handling without any LLM involvement.

In Run B (AYLA), FREQ-01 improves dramatically to 9/12 = 75.0\%, and
RTCI-01 fails across a broader variety of misattributions (R4-PM $\times 3$,
PHASE-01 $\times 1$, R3-FDIA $\times 1$, EXPECTED-09 $\times 1$,
R2-FDIA+R4-PM+CURR-01 $\times 1$), making the 0/7 failure rate less
predictable.

\textbf{Practical implication}: For deployments where RTCI-01 events are
independently captured by Snort IDS threshold alerts (which directly flag the
log spikes that characterize RTCI-01), Run A's concentrated and predictable
failure mode is operationally more manageable than AYLA's distributed
misattributions---the operator knows exactly which event class to route
around the LLM.
AYLA is the preferred choice when comprehensive, protocol-accurate attribution
across all 13 protocols is required and no external Snort-based RTCI-01
detection is available.

% ─────────────────────────────────────────────────────────────────────────────
\section{Discussion}
\label{sec:discussion}
% ─────────────────────────────────────────────────────────────────────────────

\subsection{Remaining Failure Modes}

RTCI-01's 0/7 failure rate is the dominant unsolved problem.
RTCI-01 involves current collapse correlated with Snort IDS log spikes;
however, the ensemble's SHAP attribution frequently elevates phase-angle or
current-harmonic sensors above the log channels when both co-occur.
When log channels do not surface as the primary SHAP driver, the RAG query
lacks the vocabulary to distinguish RTCI-01 from relay-specific FDIA scenarios.
A targeted fix was attempted: promoting Snort log channel drivers in the RAG
query when they are active, regardless of relative SHAP rank.
However, this approach introduced new attribution failures in other protocol
classes---Snort log activity co-occurs with multiple attack types, so
unconditionally promoting log channels over-retrieves RTCI-01 for events that
are actually relay-specific FDIA or FREQ-01.
RTCI-01 therefore remains an open failure class.
Resolving it likely requires a fundamentally different approach---such as a
dedicated log-channel detection stage that treats Snort spikes as a hard
retrieval filter, or a protocol-aware query template that separates the
log-channel signal from the sensor-type signal before retrieval.

FREQ-01 improved substantially with cross-encoder reranking (0\% $\to$ 75.0\%),
confirming that the fundamental barrier was score incomparability across retrieval
passes rather than missing protocol coverage.
The three residual FREQ-01 failures occur in samples where a relay-specific
voltage sensor dominates the primary SHAP driver, causing the named-sensor query
to outweigh the supplementary frequency-deviation query even after summation.

CURR-01's 1/4 accuracy and DEGR-02's 1/2 accuracy reflect their small sample
counts (4 and 2, respectively); additional evaluation samples would clarify whether
these are systematic failure modes or statistical noise.

The persistent PHASE-01 bias (3 false citations across 200 samples) reflects
a document-level issue: PHASE-01's protocol text contains broad language about
``voltage anomalies'' and ``sensor magnitude deviations'' that overlaps with
FDIA attack descriptions.
Tightening PHASE-01's language to emphasize its unique discriminator
(simultaneous step-change in phase angle \textit{with stable magnitudes})
would reduce its embedding similarity to voltage-magnitude queries.

\subsection{Hardware and Deployment Requirements}

\textbf{Current test hardware.} All AYLA benchmark results were produced on
an NVIDIA RTX 3060 laptop GPU (6~GB VRAM, $\sim$80~W TDP).
Pipeline throughput is approximately \textbf{3.5 minutes per event} when the
full AYLA pipeline is triggered: SHAP attribution requires approximately 30--60~s,
RAG retrieval approximately 10--20~s, and LLM generation approximately 100--120~s
(using INT4-quantized Mistral-Nemo at 4-bit precision to fit within 6~GB VRAM).
The full 200-sample benchmark required approximately 12 hours of wall-clock time.

\textbf{Key architectural note.} The LLM stage is event-driven: it is triggered
\textit{only} when a sensor snapshot exceeds the classification threshold $\tau$.
In a real grid monitoring system receiving thousands of readings per minute,
the vast majority of readings are processed at Stage~1 speed (sub-millisecond
LightGBM classification); the 3.5-minute LLM cost is incurred only for events
that are flagged as anomalies.

\textbf{Production hardware estimates.} Throughput scales with GPU memory
bandwidth, which determines LLM token generation speed:

\begin{itemize}
    \item \textbf{NVIDIA RTX 4090} (24~GB VRAM, $\sim$1,008~GB/s bandwidth,
          $\sim$3$\times$ the laptop 3060's bandwidth): estimated \textbf{45--90~s}
          per pipeline event, enabling FP16 inference instead of INT4 quantization.
          At this speed, a single workstation handles approximately
          \textbf{40--80 flagged events per hour}.

    \item \textbf{NVIDIA A100} (80~GB VRAM, $\sim$2,039~GB/s bandwidth,
          $\sim$6$\times$ improvement): estimated \textbf{20--40~s} per event,
          corresponding to approximately \textbf{90--180 events per hour}.

    \item \textbf{Multi-substation deployment}: a clustered GPU server with
          2--4$\times$ A100s provides real-time concurrent coverage across
          multiple relay systems.
\end{itemize}

\textbf{Practical conclusion.} For single-substation deployment where attack
events are rare (1--10 per day under normal threat conditions), even the current
laptop-class GPU is viable for production use.
For multi-substation control centers monitoring dozens of relay systems
simultaneously, a dedicated GPU server (RTX~4090 or A100-class) is recommended.

\subsection{Relation to IHGR-RAG}

IHGR-RAG \cite{ihgrrag2025}, developed for power equipment condition assessment,
employs a pipeline strikingly similar to AYLA: hierarchical retrieval, global context,
LLM generation, and CoT-based text alignment.
The AYLA pipeline directly adopts the cross-encoder reranking and CoT paradigm
from IHGR-RAG \cite{ihgrrag2025} and extends it to the SCADA relay security domain
with SHAP-attribution-guided query construction.
In doing so, AYLA validates and extends the IHGR-RAG architecture in a novel domain,
demonstrating that CoT+cross-encoder reranking generalizes from power equipment
condition assessment to cyber-physical attack attribution.
The primary remaining methodological difference is that IHGR-RAG employs hierarchical
retrieval with global context injection, while AYLA uses SHAP-attribution-guided
primary queries followed by type-pattern secondary passes---a domain-specific
retrieval design motivated by the relay-sensor vocabulary of ICS protocols.

\subsection{Real-World Deployment Considerations}

In a production grid operations center, AYLA would operate as a second-tier
alerting layer behind the existing SCADA IDS.
The 14.7\% false alert rate at $\tau{=}0.560$ is competitive with industrial
thresholds for grid anomaly detection and can be adjusted via the threshold
sweep (Table~\ref{tab:sweep}) without retraining.
The local LLM deployment (Mistral-Nemo via Ollama) is critical for environments
where SCADA telemetry is classified and cannot be transmitted to external API endpoints.

Unlike Claroty, Dragos, and Nozomi Networks---which surface alerts but provide
no automated, grounded remediation guidance---AYLA delivers step-by-step
protocol-grounded remediation advice for each flagged event.
This directly reduces the manual documentation lookup burden that operators
currently face during live incidents.
NERC CIP compliance requires incident logging but not automated response;
AYLA's explainable attribution provides an audit trail that supports compliance
reporting while simultaneously assisting operators.

The cascaded per-scenario classifier (Section~\ref{subsec:cascade}) provides a
complementary fast path.
By routing the majority of alerts through the cascade (sub-millisecond) and
reserving the full AYLA pipeline for ambiguous events, mean response latency
drops substantially without sacrificing attribution accuracy on high-confidence
cases.

% ─────────────────────────────────────────────────────────────────────────────
\section{Conclusion}
\label{sec:conclusion}
% ─────────────────────────────────────────────────────────────────────────────

This paper presented Power Shield, implementing an end-to-end
Automated Yield-Loss Alerting (AYLA) pipeline for power grid attack detection
and LLM-guided remediation.
Three principal results establish the pipeline's feasibility.

\textbf{First, classifier feasibility}: the SHAP-guided LightGBM/XGBoost stacked
ensemble achieves AUC-ROC 0.9665, MCC 0.8020, and attack recall 94.6\% at
$\tau{=}0.560$, surpassing the Hink et al.\ \cite{hink2014ml} combined-dataset
baseline (90.56\%) on the full heterogeneous MSU/ORNL dataset.

\textbf{Second, LLM grounding confirmed by ablation}: AYLA achieves 100\%
RAG retrieval accuracy across all 13 protocol test cases and 90.5\% correct
protocol citation (181/200, 95\% CI [85.6\%, 93.8\%]).
A no-context ablation---identical prompts with retrieved documents withheld---
drops attribution accuracy to 16.5\% (33/200, 95\% CI [12.0\%, 22.3\%]),
a 74.0~pp decline with non-overlapping confidence intervals that directly
confirms the LLM is grounding its responses in the retrieved protocol manuals.
A $\chi^2$ test ($\chi^2 = 1931.42$, $p < 10^{-10}$) further confirms
full-pipeline performance far exceeds the random-selection baseline.

\textbf{Third, comparison with deployed systems}: unlike current commercial ICS
security platforms (Claroty, Dragos, Nozomi Networks), which provide passive
network monitoring and alert surfacing but no automated remediation guidance,
AYLA delivers step-by-step, protocol-grounded remediation advice for each
flagged event---a capability absent from all known deployed ICS security systems.

An architecture comparison between the original single-LightGBM system (Run A,
84.5\% LLM accuracy) and AYLA (Run B, 90.5\%) shows a 6.0~pp improvement
attributable to the ensemble classifier and cross-encoder reranking.
The comparison also reveals a shift in failure character: Run A has concentrated,
predictable failures (RTCI-01 always misattributed to R4-PM; FREQ-01 always
fails), while AYLA has distributed failures (RTCI-01 is the only systematic
failure mode).

The principal remaining limitation is RTCI-01 attribution (0/7 = 0\% in AYLA),
where phase-angle SHAP drivers suppress the log-channel vocabulary needed to
distinguish RTCI-01 from relay-specific FDIA scenarios.
A log-channel promotion fix was attempted but introduced new failures in other
protocol classes, confirming that RTCI-01 resolution requires a fundamentally
different retrieval strategy rather than a SHAP-rank override.
Future work will focus on dedicated log-channel detection stages and integration
of an online retraining loop to adapt the classifier to concept drift in
adversarial grid telemetry.

% ─────────────────────────────────────────────────────────────────────────────
\bibliographystyle{unsrt}
\bibliography{references}
% ─────────────────────────────────────────────────────────────────────────────

\end{document}
